\section{Implementation}
We implement the model in System W, an extensible .NET compiler. The contract owner is a production module shared by selection, compilation observers, and experiment runners. A fact state exposes \texttt{valid}, \texttt{invalid}, and \texttt{unknown}; an effect contract contains normalized \texttt{requires}, \texttt{produces}, \texttt{preserves}, and \texttt{invalidates} sets; and each \texttt{VerificationObligation} stores fact, rule, canonical owner, creation boundary, and first eligible boundary.

The production observer receives actual bytecode, AIR, optimized AIR, and backend-input artifacts. The backend-input hook is part of the normal execution builder and runs immediately before backend compilation. A compilation-scoped lifecycle record is keyed by the identity of the immutable compilation input. It contains the last observed boundary, current fact state, and pending obligations. Boundary order must be strictly increasing; state is discarded on completion or failure so a failed attempt cannot contaminate a retry.

Stage seeds add only facts not already marked invalid. Effect validation then follows selected module occurrence order rather than sorting by module name. It rejects missing requirements, invalid preservation claims, repeated ambiguous occurrences, and invalidations without a canonical verifier route. A route descriptor records its earliest executable boundary. The created deadline is the later of creation and route availability, so an optimized-AIR invalidation can remain pending until backend input. The scheduler uses deterministic rule/fact ordering and rejects unknown policies, unknown demand queries, route conflicts, owner mismatches, and missed deadlines.

P0, P1, P1D, P2, and P3 share the same structural and semantic verifier implementations. P1D adds only an execution-scoped demand set; it does not alter the artifact or oracle. P2 schedules due obligation routes. P3 schedules every route exposed at the boundary. A route may own several facts, so selectivity is at the canonical-route granularity: one invocation can re-establish multiple facts. We report verifier invocations rather than presenting fact count as checker cost.

\subsection{Admissibility and boundary procedure}
Contract admission is separated from policy execution. Selection first normalizes fact identifiers, effect sets, route identifiers, owner identities, and earliest executable boundaries. It then resolves every selected fact to at most one canonical verifier owner and rejects duplicate owners, an unknown boundary, or an invalidation with no executable route. This preflight prevents a later policy choice from changing which checker is authoritative. It also keeps P1D, P2, and P3 comparable: they differ in scheduling, not in verifier implementation or route resolution.

At each observed boundary the implementation performs four ordered steps. First, it merges the stage seed into the retained state without overwriting a prior invalidation, then validates requirements and preservation claims. Second, it applies produced and invalidated effects in selected occurrence order; each invalidation creates or retains an obligation whose lineage records the triggering boundary. Third, it computes the execution set: P1D chooses queried invalid facts whose routes are executable, P2 chooses obligations whose first eligible boundary is current, and P3 chooses every canonical route exposed there. Fourth, it executes routes in deterministic rule/fact order. A successful invocation marks its owned scheduled facts valid and removes matching due obligations. Rejection, an unknown result, an owner mismatch, an unresolved route, an out-of-order boundary, an overdue obligation, or pending P2/P3 work at the final modeled boundary stops compilation.

Structural and semantic establishment are intentionally separate. AIR receives one semantic baseline when it first becomes observable. If P1 later records an invalidation, that baseline is not treated as a discharge event. Conversely, a P3 invocation with no obligation can establish the route's owned facts at its current boundary. This prevents mandatory baseline checks from making P1 appear to enforce obligation discharge.

The implementation deliberately does not infer a route from a nearby component name or from registration order. Such inference would make correctness depend on composition accidents invisible to the effect contract. Repeated occurrences are accepted only when their normalized effects and ownership are unambiguous. This is why occurrence order, owner identity, missing-route handling, conflict rejection, and deferred-boundary handling appear as production invariants and targeted tests.

\subsection{Implementation-conformance checks}
The conformance suite checks the actual orchestration path, not only the standalone scheduler. It asserts the observer order
\[
\begin{aligned}
\text{bytecode}&<\text{AIR}<\text{optimized AIR}\\
&<\text{backend input}<\text{backend compile}.
\end{aligned}
\]
A deferred case invalidates \texttt{backend.input-verified} at optimized AIR. P2 creates an obligation with creation boundary optimized AIR and first eligible boundary backend input, carries it in the compilation-scoped state, and invokes the backend-input owner before compilation. Separate tests verify carry-forward when no new effect reproduces the obligation, overdue rejection, route/owner conflicts, and state cleanup after failure. These tests establish code/model correspondence for lifecycle mechanics; they do not validate the semantic truth of arbitrary extension contracts.

\paragraph{System W language path.}
The source-to-result runner executes the normal dialect engine, optimizer, production observer, and CIL/interpreter backends in fresh processes. Research policy switches remain internal/friend APIs and do not expand the shipped facade. P3 remains the strongest checking profile, while the frozen historical experiment runner continues to count only boundaries represented by its versioned protocol.

\paragraph{Language T public-SDK path.}
A separate small language package uses only the public extension SDK. Its shape/layout facts and verifier route differ from System W's arithmetic/capability relations, while the scheduler and contract representations are shared. We do not call the package independently authored: its cases were created during this study.

\paragraph{Baseline defect repair.}
Historical source-to-result case P07 mixed double parameters with an integer literal and failed under every policy before a result. We retained the historical receipt, repaired the independent numeric-dispatch defect with an explicit finite promotion matrix, added CIL/interpreter regression coverage, versioned the corpus, and reran the unchanged source and oracle. P07 now produces 13 under every policy and is counted as a valid control only in the new protocol.

The exact pre-hardening master contract was 1,597 passing tests. New test counts are reported only after mechanical discovery on the final branch head; the historical count is not rewritten.
